\section{Einleitung}
\label{chap:einleitung}

\subsection{Motivation}
\label{sec:motivation}

Die zuverlässige Messung hoher Ströme in Niederspannungsschaltanlagen wird im Kontext der Energiewende und der zunehmenden Dezentralisierung der Energieversorgung immer wichtiger \cite{pdf-Zellularer_Ansatz-VDE}. Schaltanlagen bilden in industriellen Infrastrukturen zentrale Verteil- und Knotenpunkte. Entsprechend werden Strommesswerte nicht nur für Schutzfunktionen benötigt, sondern zunehmend auch für Abrechnungszwecke, ein strukturiertes Energiemanagement sowie zur Unterstützung der Netzstabilität \cite{norm-DIN_ISO_50001-Energiemanagement-DIN}. Damit steigen die Anforderungen an Genauigkeit, Reproduzierbarkeit und Langzeitstabilität – besonders im Hochstrombereich.

In der Praxis treffen diese Anforderungen auf konstruktive und wirtschaftliche Randbedingungen. Hersteller wie die \textit{Rolf~Janssen~Elektrotechnische~Werke~GmbH} müssen Schaltanlagen häufig sehr kompakt auslegen, da technische Betriebsräume nur begrenzt Platz bieten \cite[Kap.~1.4]{pdf-Planung_Energieverteilung-Siemens}. Kompakte Bauweisen erhöhen jedoch die Packungsdichte: Sammelschienen mit mehreren Kiloampere verlaufen oft nahe an Messstromwandlern.

Genau daraus entsteht der Zielkonflikt dieser Arbeit: Die Anwendung verlangt präzise Strommessung, gleichzeitig verschlechtern magnetische Streufelder benachbarter Leiter die Messbedingungen im Wandler. Hochspezialisierte Komponenten sind nicht automatisch die beste Lösung, da sie Kosten und Beschaffungsrisiken erhöhen können. Ziel ist daher, Messgüte und Wirtschaftlichkeit so zu kombinieren, dass eine serientaugliche Lösung für eine neue Schaltanlagengeneration entsteht.

\subsection{Problemstellung}
\label{sec:problemstellung}

Moderne Niederspannungsschaltanlagen vereinen hohe Stromtragfähigkeit mit geringer Baugröße. Dadurch befinden sich stromführende Leiter und Messstromwandler in engem räumlichem Verbund. Messstromwandler werden dabei magnetischen Fremdfeldern ausgesetzt, die von benachbarten Leitern ausgehen und die Messerfassung beeinflussen können \cite{pdf-Fremdfeldstoerungen_Stromwandler-Redur}. In Drehstromsystemen ist dieser Effekt besonders relevant: Die Phasenleiter L1, L2 und L3 sind typischerweise parallel geführt oder räumlich dicht gruppiert. Die zeitlich veränderlichen Magnetfelder der Nachbarphasen können in den Wandler einkoppeln und den magnetischen Arbeitspunkt im Kern verschieben. Daraus resultieren Messabweichungen, die mit steigender Stromstärke deutlich zunehmen können.

\begin{figure}[H]
    \centering
    \begin{subfigure}[b]{0.45\textwidth}
        \centering
        \includegraphics[width=\textwidth]{03_Ressourcen/Bilder/schaltschrank_front.jpg}
        \caption{Frontansicht}
        \label{fig:schaltschrank_front}
    \end{subfigure}
    \hfill
    \begin{subfigure}[b]{0.45\textwidth}
        \centering
        \includegraphics[width=\textwidth]{03_Ressourcen/Bilder/schaltschrank_back.jpg}
        \caption{Rückansicht}
        \label{fig:schaltschrank_back}
    \end{subfigure}
    \caption{Gesamtansicht des untersuchten Schaltschranksystems}
    \label{fig:schaltschrank_gesamt}
\end{figure}

Die Messabweichungen wirken sich unmittelbar auf die Zuverlässigkeit von Schutzkonzepten, Betriebsdatenerfassung und Verrechnungsprozessen aus. Für den betrachteten Aufbau zeigen die Messdaten eine deutliche Fremdfeldbeeinflussung ab \SI{2500}{A} (vgl. Tabelle~\ref{tab:messergebnisse_celsa10030_2500A_alle}). Im Hochstrombereich können sich Abweichungen im Prozentbereich einstellen. Zusätzlich tritt die Beeinflussung im Drehstromsystem nicht zwingend symmetrisch auf: Abhängig von Leiterführung, Abständen und Einbausituation kann eine Phase stärker betroffen sein als die anderen. In der untersuchten Schaltanlage zeigt sich dies insbesondere auf Phase L2. L2 wird daher als maßgeblicher Worst-Case für die Bewertung herangezogen.

Am Markt existieren verschiedene Gegenmaßnahmen, z.\,B. kompensierte Wandler oder zusätzliche Abschirmungen. Diese Ansätze sind jedoch mit Zielkonflikten verbunden: Kompensierte Ausführungen erhöhen häufig Kosten und Integrationsaufwand, während Abschirmmaßnahmen stark von Einbaulage, Toleranzen und verfügbarem Bauraum abhängen. Damit stellt sich die Frage, ob sich durch eine gezielte System- und Geometrieoptimierung auch mit wirtschaftlichen Standardwandlern eine normkonforme Messgüte unter realen Hochstrombedingungen erreichen lässt. Genau diese Frage adressiert die vorliegende Arbeit.

\subsection{Zielsetzung und Leitfrage}
\label{sec:zielsetzung}

Ziel der Arbeit ist es, für die Neuentwicklung einer kompakten Niederspannungsschaltanlage eine Lösung zur Strommessung zu bewerten, die technisch belastbar, wirtschaftlich vertretbar und konstruktiv integrierbar ist. Untersucht wird die Messung im Drehstromsystem (L1, L2, L3). Aufgrund der ausgeprägten Beeinflussung wird L2 als kritischer Bewertungsfall genutzt.

Die zentrale Leitfrage lautet:
\emph{Welche Maßnahmen an Wandler, Schutzkonzept und Leiter-/Wandlergeometrie stellen sicher, dass die Strommessung auf L2 auch unter Fremdfeldeinfluss im Hochstrombereich die geforderte Messgenauigkeit einhält?}

Zur Beantwortung werden drei Lösungsrichtungen betrachtet, die unterschiedliche Ebenen adressieren:

\begin{itemize}
    \item \textbf{Einsatz kompensierter Wandler (Technologieansatz):}
    Untersuchung der Fremdfeldrobustheit kompensierter Ausführungen und Einordnung des technischen Nutzens im Verhältnis zu Kosten, Bauraum und Integrationsaufwand.
          
    \item \textbf{Fremdfeld-Protektion, FFP (konstruktiver Schutzansatz):}
    Bewertung, inwieweit konstruktive Schutzmaßnahmen die Einkopplung externer Felder reduzieren und wie empfindlich die Wirksamkeit gegenüber Montage- und Toleranzeinflüssen ist.
          
    \item \textbf{Geometrische Optimierung mittels Dreiecksanordnung mit Standardwandler (Systemansatz):}
    Prüfung, ob eine angepasste Leiter-/Wandleranordnung die Feldüberlagerung im Drehstromsystem so entschärft, dass Standardwandler normnahe Ergebnisse liefern.
\end{itemize}

Als Bewertungsmaßstäbe dienen die Messabweichung \gls{sym:epsilon}, die Stabilität unter Fremdfeldeinfluss (insbesondere auf L2) sowie eine Kosten-/Bauraum-Einordnung im Hinblick auf die Serienumsetzung.

\subsection{Vorgehensweise und Aufbau der Arbeit}
\label{sec:vorgehensweise}

Die Arbeit umfasst einen theoretischen Grundlagenteil, die Beschreibung des Prüfaufbaus und die Auswertung der Messergebnisse. Kapitel~\ref{chap:theorie} erläutert die physikalischen Grundlagen der Stromwandlung sowie die Mechanismen magnetischer Induktion und Kernsättigung. Zusätzlich werden relevante Normen und Genauigkeitsklassen als Bewertungsrahmen eingeführt.

Kapitel~\ref{chap:methodik} beschreibt den Hochstromprüfstand, die Messkette und deren Auslegung. Darüber hinaus werden die untersuchten Wandlerkonzepte sowie die betrachteten Leitergeometrien und Anordnungen im Drehstromsystem definiert.

Kapitel~\ref{sec:auswertung_diskussion} wertet die Messreihen im Bereich von \SI{2000}{A} bis \SI{4000}{A} aus. Die Ergebnisse werden phasenbezogen analysiert, mit Schwerpunkt auf \gls{sym:epsilon}, Fremdfeldrobustheit und der Bewertung des Worst-Case L2. Ergänzend erfolgt eine wirtschaftliche Einordnung, um technische Wirkung und Umsetzbarkeit gegenüberzustellen.

Kapitel~\ref{chap:zusammenfassung_ergebnisse} fasst die Ergebnisse zusammen, leitet eine Handlungsempfehlung für die Neuentwicklung ab und formuliert konkrete nächste Schritte zur Überführung in eine serienfähige Konstruktion.
